% Source-derived chapter 01: Introduction
% Original source: birational-section-conjecture-strong-birationality
\chapter{Introduction}
\label{birational-section-conjecture-strong-birationality-chapter-01}
\source{birational-section-conjecture-strong-birationality}

\begin{remark}[Source section]
\label{birational-section-conjecture-strong-birationality-chapter-01-source}
\source{birational-section-conjecture-strong-birationality}
\source{birational-section-conjecture-strong-birationality}
This chapter reproduces the corresponding section of the retrieved
LaTeX source, including its definitions, statements, examples, and
proofs. It has not yet been translated into Lean.
\notready
\end{remark}

% The source section title was: Introduction
Given $X$ a geometrically connected, smooth curve over a field $k$ with absolute Galois group $\gal(k)$, let $\mc{S}_{X/k}$ be the set of sections of $\pi_{1}(X)\to\gal(k)$ modulo the action of $\pi_{1}(X_{\bar{k}})$ by conjugation, these are usually called Galois sections. A Galois section is \emph{geometric} if it is associated with a $k$-rational point of $X$ and \emph{cuspidal} if it is associated with a $k$-rational point in the boundary $\bar{X}\setminus X$, where $\bar{X}$ is the smooth completion of $X$. Grothendieck made the following conjecture, nowadays called the \emph{section conjecture}, in a letter to Faltings \cite{gro97}.

\begin{SC}
\source{birational-section-conjecture-strong-birationality}
\notready
    For every smooth, geometrically connected hyperbolic curve $X$ over a field $k$ finitely generated over $\Q$, every Galois section of $X$ is either geometric or cuspidal.
\end{SC}

A birational version of the conjecture, called the \emph{birational section conjecture}, has been studied as it seems more approachable \cite{koe05} \cite{sti15} \cite{saty21}. We can define a set $\mc{S}_{k(X)/k}$ of sections of $\gal(k(X))\to\gal(k)$ analogously to the above, its elements are called birational Galois sections. A birational section is cuspidal if it comes from a rational point of $\bar{X}$.

\begin{BSC}
\source{birational-section-conjecture-strong-birationality}
\notready
    For every smooth, geometrically connected curve $X$ over a field $k$ finitely generated over $\Q$, every birational Galois section of $X$ is cuspidal.
\end{BSC}

Recall that a Galois section of $\mc{S}_{X/k}$ is \emph{birationally liftable} if it's in the image of $\mc{S}_{k(X)/k}\to\mc{S}_{X/k}$ \cite[\S 2.1.1]{sti15} (this is slightly different from the definition given in \cite[Definition 266]{sti13}). Here is an alternative formulation of the birational section conjecture, see \autoref{bireq}.

%Clearly, the birational section conjecture implies that every birationally liftable section of $X$ is either geometric or cuspidal. On the other hand, if $s\in \mc{S}_{k(X)/k}$ is a birational section whose image in $\mc{S}_{U/k}$ is geometric or cuspidal for every open subset $U\subset X$, then $s$ is cuspidal since $\mc{S}_{k(X)/k}\simeq\projlim_{U}\mc{S}_{U/k}$ \cite[Lemma 259]{sti13}. Because of this, the birational section conjecture is equivalent to the following.

\begin{BSC2}
\source{birational-section-conjecture-strong-birationality}
\notready
    For every smooth, geometrically connected curve $X$ over a field $k$ finitely generated over $\Q$, every birationally liftable Galois section of $X$ is either geometric or cuspidal.
\end{BSC2}

% Hence, the birational section conjecture for all curves over $k$ is equivalent to the statement that every birationally liftable section of every curve over $k$ is geometric or cuspidal.

%The birational section conjecture states that birationally liftable sections are geometric or cuspidal.

In this note we show that, if we strengthen the birationality assumption, we can obtain the desired result. Namely, we show that if $t$ is an indeterminate and $s$ is a Galois section such that the base change $s_{k(t)}$ is birationally liftable, then $s$ is geometric or cuspidal. 

\subsection{Known results}

J. Koenigsmann \cite{koe05} proved that the birational section conjecture holds over finite extensions of $\Q_{p}$. Clearly, one would like to pass from local fields to number fields. Moreover, M. Sa\"idi and M. Tyler \cite{saty21} have proved that the birational section conjecture for number fields implies it for finitely generated extensions of $\Q$.

J. Stix \cite{sti15} obtained partial results about the passage from local fields to number fields, let us describe them. Fix $X$ a smooth, projective curve over a number field $k$ and $s\in\mc{S}_{k(X)/k}$ a Galois section of $k(X)$. Using Koenigsmann's results, $s$ induces a point $x_{\nu}\in X(k_{\nu})$ for every place $\nu$.

Given an open subset $U\s X$, choose some spreading out $\tilde{U}\to\spec \mf{o}_{k,S}$ of $U$, where $\mf{o}_{k}$ is the ring of integers of $k$ and $S$ is some finite set of places. Let $N_{U}$ be the set of places of $k$ such that $x_{\nu}$ is \emph{not} integral: we have that $N_{U}$ depends on the choice of $\tilde{U}$ only up to a finite number of places, hence the Dirichlet density of $N_{U}$ is well defined.

Stix first proves that, if $k$ is a totally real or imaginary quadratic number field, then $N_{U}$ is infinite for some open subset $U\s X$ \cite[Theorem A]{sti15}. Secondly he proves that, if $N_{U}$ has strictly positive Dirichlet density for some open subset $U\s X$, then $s$ comes from a $k$-rational point of $X\setminus U$ \cite[Theorem B]{sti15}. In order to prove the birational section conjecture for totally real or imaginary quadratic number fields, it is then sufficient to bridge the gap between Stix's two results.

\subsection{Our main theorem}

We study the passage from local fields to number fields, too, but we use a different approach. We strengthen the birationality assumption: under this strengthened hypothesis, we obtain a complete result.

\begin{definition*}
\source{birational-section-conjecture-strong-birationality}
\notready
    Let $X$ be a curve over a field $k$, $s\in\mc{S}_{X/k}$ a Galois section and $t$ an indeterminate; consider the base change $s_{k(t)}\in\mc{S}_{X_{k(t)}/k(t)}$ of $s$ to $k(t)$ \cite[Definition 27]{sti13}. We say that $s$ is \emph{$t$-birationally liftable} if $s_{k(t)}$ is birationally liftable.
\end{definition*}

\begin{TA}
\source{birational-section-conjecture-strong-birationality}
\notready
    \hypertarget{TA}{Let} $X$ be a smooth curve over a field $k$ finitely generated over $\Q$. A Galois section of $X$ is geometric or cuspidal if and only if it is $t$-birationally liftable.
\end{TA}

%~ While the proof of \hyperlink{TA}{Theorem A} is quite technical, the underlying idea is simple. 
A common theme in the study of the section conjecture is the attempt to replicate with Galois sections the geometric constructions we make with points, see e.g. \cite[Chapter 3]{sti13}. The idea underlying the proof of \hyperlink{TA}{Theorem A} is, given $s$ a Galois section of $\P^{1}\setminus\{0,1,\infty\}$, to define ``the polynomial $t-s$'' so that, if $s$ is geometric associated to a $k$-rational point $x \in k\setminus\{0,1\}$, then ``$t-s$'' coincides with the polynomial $t-x$. If $s$ is $t$-birationally liftable then we are able to define this ``polynomial'' by using a birational lifting of $s_{k(t)}$ to produce an element of $\H^{1}(k(t),\hz(1))=\hat{k(t)^{*}}$.

 %~ ``a polynomial with coefficients in non-cuspidal Galois sections of $\P^{1}\setminus\{0,1,\infty\}$'', which conjecturally should correspond to $k\setminus\{0,1\}$. Given $s$ a Galois section of $\P^{1}\setminus\{0,1,\infty\}$, if $s$ is $t$-birational then we are able to define the ``the polynomial $t-s$'' by using a birational lifting of $s_{k(t)}$ to produce an element of $\H^{1}(k(t),\hz(1))=\hat{k(t)^{*}}$ which, if $s$ is geometric associated to a $k$-rational point $x$, coincides with the image of $t-x$ with respect to the natural map $k(t)^{*}\to\hat{k(t)^{*}}$.

This turns out to be sufficient: if we are over a number field $k$ and, for a place $\nu$, the Galois section $s_{k_{\nu}}$ is associated with a point $x_{\nu}\in k_{\nu}\setminus\{0,1\}$, the fact that the ``polynomial $t-s=t-x_{\nu}$'' is defined over $k$ forces $x_{\nu}$ to be $k$-rational and not to depend on $\nu$.

%~ Translating this idea into a proof requires a number of technical tools. One of them is understanding specialization for \emph{ramified} Galois sections.  specialization is done classically only for \emph{unramified} sections, see \cite[Chapter 8]{sti13}. We develop such a theory in \autoref{ramspe}. 
%~ We use a non-standard valuative criterion for proper morphisms of Deligne-Mumford stacks: this allows us to specialize Galois sections just as we specialize points of schemes using the valuative criterion for properness. 



\subsection{Consequences for the section conjecture}

One of the main reasons for studying the birational section conjecture is to reduce the section conjecture to a \emph{lifting}, or \emph{cuspidalization}, problem.

The section conjecture easily implies the following statement, which is sometimes called \emph{cuspidalization conjecture} in the literature.

\begin{CC}
\source{birational-section-conjecture-strong-birationality}
\notready
    For every smooth, geometrically connected hyperbolic curve $X$ over a field $k$ finitely generated over $\Q$ and every non-empty open subset $U\subset X$, the map $\mc{S}_{U/k}\to\mc{S}_{X/k}$ is surjective.
\end{CC}

By a limit argument, the cuspidalization conjecture holds if and only if every Galois section of every hyperbolic curve over $k$ is birationally liftable, see \autoref{cuspbl}.

H. Esnault and P.H. Hai showed that the cuspidalization conjecture reduces the section conjecture to the case of $\P^{1}\setminus\{0,1,\infty\}$ \cite[Proposition 7.9]{eh08}. For $\P^{1}\setminus\{0,1,\infty\}$ one may hope to solve the conjecture with explicit computations, such as the $n$-nilpotents obstructions introduced by J. Ellenberg and K. Wickelgren \cite{wic12}. Moreover, the cuspidalization conjecture reduces the section conjecture to the birational section conjecture. Neither the case of $\P^{1}\setminus\{0,1,\infty\}$ nor the birational section conjecture is needed: the section conjecture is equivalent to the cuspidalization conjecture.


\begin{TB}
\source{birational-section-conjecture-strong-birationality}
\notready
	\hypertarget{TB}{Let} $X$ be a hyperbolic curve over a field $k$ finitely generated over $\Q$. The following are equivalent.
	\begin{itemize}
		\item For every finitely generated extension $K/k$, every Galois section of $X_{K}$ is either geometric or cuspidal.
		\item For every finitely generated extension $K/k$, every Galois section of $X_{K}$ is birationally liftable.
	\end{itemize}
        As a consequence, the section conjecture is equivalent to the cuspidalization conjecture.
\end{TB}



\subsection{Consequences for the birational section conjecture}

It is well known that it is enough to prove the birational section conjecture for $\P^{1}\setminus\{0,1,\infty\}$. Thanks to \hyperlink{TA}{Theorem A}, it is then enough to show that for every birationally liftable section $s$ of $\P^{1}\setminus\{0,1,\infty\}$ over a number field $k$ the base change $s_{k(t)}$ lifts to $\gal(k(t)(\P^{1}))$. We prove that it is enough to find a much simpler lifting, and we manage to reduce the existence of said lifting to a problem of ``interpolation of Galois sections'', i.e. finding a Galois section with prescribed specializations of some curve over $k(t)$. Let us explain this.

We say that a morphism $X\to V$ is a \emph{family of curves} if it has the form $\bar{X}\setminus D\to V$ where $\bar{X}\to V$ is smooth, proper, with geometrically connected fibers of dimension $1$, and $D\subset\bar{X}$ is a divisor finite étale over $V$. Assume that $V$ is an affine curve over a field of characteristic $0$ and let $F$ be a geometric fiber, then the sequence
\[1 \to \pi_{1}(F)\to\pi_{1}(X)\to\pi_{1}(V)\to 1\]
is exact, see \autoref{families} (we stress that we are assuming $V$ affine, otherwise this is false).

We may define a space $\mc{S}_{X/V}$ of ``global'' sections $\pi_{1}(V)\to\pi_{1}(X)$ modulo conjugation by $\pi_{1}(F)$. Let us call these $\pi_{1}$ sections, in analogy with Galois sections. If $V$ is an affine open subset of $\P^{1}$ over a number field $k$, a $\pi_{1}$ section $r$ of $X\to V$ is uniquely determined by its specializations $r_{v}\in\mc{S}_{X_{v}/k}$, $v\in V(k)$, see \autoref{hilbsp}.
%~ : this follows from the fact that the subgroup generated by the sections $\gal(k)\to\pi_{1}(V)$ associated with rational points of $V$ is dense by Hilbert's irreducibility theorem, see \autoref{hilbgen}.

If $z\in\mc{S}_{k(\P_{1})/k}$ is a birational Galois section of $\P^{1}$, for every open subset $U\subset\P^{1}$ denote by $\iota_{U}(z)$ the image of $z$ in $\mc{S}_{U/k}$. Let $\Delta\subset\P^{1}\times\P^{1}$ be the diagonal.

%Assume that $z$ is cuspidal associated with a point $p\in\P^{1}(k)$ and let $U\subset\P^{1}$ be an affine open subset, write $V=U\setminus\{p\}$. Using the fact that $\iota_{U}(z)$ is geometric or cuspidal, we show that $\iota_{U}(z)\times V\in\mc{S}_{U\times V/V}$ lifts to a ``horizontal, geometric or cuspidal'' $\pi_{1}$ section $r\in \mc{S}_{(U\times V\setminus\Delta)/V}$ such that, for every $v\in V$, the specialization of $r$ at $v$ is $\iota_{U\setminus \{v\}}(z)$. Moreover, we prove that the converse holds: if we are able to ``interpolate'' the sections $\iota_{U\setminus \{v\}}(z)$ for every open subset $U\subset \P^{1}$ and some $V\subset U$, then $z$ is cuspidal. We thus obtain the following.

%~ Let $V\subset U$ be open subsets of $\P^{1}$ and $\Delta:V\to U\times V$ be the diagonal. Given a birational section $s\in\mc{S}_{U/k}$ with lifting $z\in\mc{S}_{k(\P^{1})/k}$, we may extend $s$ to a ``horizontal'' $\pi_{1}$ section $s\times V$ of $U\times V\to V$. If $z$ is cuspidal associated with a point $p\in\P^{1}(k)$, we may lift $z_{k(t)}$ to a cuspidal section of $\spec k(U)\times V\setminus\Delta$ and define $r_{U}$ as the image of said lifting in $U\times V\setminus\Delta \to V$. For every $v\in V\setminus\{p\}$, the specialization of $r_{U}$ at $v$ is geometric associated with $p$, and this coincides with $z|_{U\setminus \{v\}}$. Because of this, it is natural to try to construct $r$ by gluing the sections $z|_{U\setminus \{v\}}$.

%~ We prove that $s$ comes from geometry if and only if $s\times V$ lifts to a $\pi_{1}$ section $r$ of $U\times V\setminus\Delta \to V$ for some open subset $V\s U$.

%~ Moreover, we have strong clues about this hypothetical lifting $r$. 

\begin{TC}
\source{birational-section-conjecture-strong-birationality}
\notready
	\hypertarget{TC}{The} following are equivalent.
	\begin{itemize}
		\item The birational section conjecture holds.
		%~ \item For every number field $k$, every open subset $U\s\P^{1}$ and every birational section $s\in\mc{S}_{U/k}$, there exists an open subset $V\s U$ such that the extension $s\times V$ of $s$ to $U\times V\to V$ lifts to $U\times V\setminus\Delta\to V$.
		\item For every number field $k$, every section $z\in\mc{S}_{k(\P_{1})/k}$ and every open subset $U\s\P^{1}$, there exists an open subset $V\s U$ and a $\pi_{1}$ section $r$ of $U\times V\setminus\Delta \to V$ such that the specialization $r_{v}$ is equal to $\iota_{U\setminus\{v\}}(z)$ for every $v\in V(k)$.
		%~ \item For every number field $k$, every open subset $U\s\P^{1}$ and every birational section $s\in\mc{S}_{U/k}$ with lifting $z\in\mc{S}_{k(\P_{1})/k}$, there exists an open subset $V\s U$ and a representative $\gal(k)\to\pi_{1}(U\setminus\{v\})$ of $z|_{U\setminus\{v\}}$ for every $v\in V(k)$ such that the kernel of $\ast_{v}\gal(k)\to\pi_{1}(V)$ is mapped to the identity by the induced homomorphism $\ast_{v}\gal(k)\to\pi_{1}(U\times V\setminus\Delta)$, where $\ast$ denotes the free product in the category of pro-finite groups. 
	\end{itemize}
\end{TC}

Since a $\pi_{1}$ section is uniquely determined by its specializations, the section $r$ in the statement of Theorem C is a lifting to $U\times V \setminus \Delta$ of the base change $(\iota_{U}(z))_{V}\in\mc{S}_{U\times V/V}$.

%~ In short, Theorem C reduces the birational section conjecture to proving that if $z$ is a Galois section of $\gal(k(\P^{1}))\to\gal(k)$ and $U\s\P^{1}$ is an open subset, it is possible to find an open subset $V\s U$ and "compatible" representatives of the sections $z|_{U\setminus v}$, $v\in V(k)$.


\subsection*{Acknowledgements}

I would like to thank an anonymous referee for providing a very large number of comments. This note has become substantially better thanks to his or her work.

\subsection*{Notation and conventions}

Throughout the article, it is tacitly assumed that schemes do not have points of positive characteristic. In particular, fields are of characteristic $0$. Curves are smooth and geometrically connected. If $X$ is a curve, we denote by $\bar{X}$ its smooth completion. The letter $R$ always denotes a DVR with fraction field $K$ and residue field $k$. If $A$ is an abelian group, we denote by $\hat{A}$ the projective limit $\projlim_{n}A/nA$. We write b.l. (resp. $t$-b.l.) as an abbreviation for birationally liftable (resp. $t$-birationally liftable).



%~ If $A$ is an abelian group and $\ell$ is a prime, we write $\hat{A}[\ell^{\infty}]$ for the pro-$\ell$ completion $\projlim_{n}A/\ell^{n}A$ and $\hat{A}$ for $\projlim_{n}A/nA$. The group $\hat{A}[\ell^{\infty}]$ is the projective limit of the $\ell$-primary quotients of $A$, and should not be confused with the $\ell$-primary torsion subgroup of $\hat{A}$. Since we will use the pro-$\ell$ completion $\hat{k}^{*}[\ell^{\infty}]$ of the group of non-zero elements of a field $k$, we do not use the more standard notation $A_{\ell}$ to avoid confusion with the completion of a field at a place. There is a natural isomorphism $\hat{A}\simeq\prod_{\ell}\hat{A}[\ell^{\infty}]$.



%~ The isomorphism classes of points of the fibers are in natural bijection with the space of Galois sections of the étale fundamental group of the fiber. In short, $\Pi_{X/S}\to\spec S$ is a compact and coherent way of packing the spaces of Galois sections of the fibers.

%~ In \cite[\S 1]{scfg} we generalized the étale fundamental gerbe to a relative setting: this gives a suitable framework to specialize Galois sections using fundamental gerbes.

%~ Let $S,X$ be noetherian, normal, connected schemes and let $X\to S$ be a geometric fibration with connected fibers (see \cite[Definition 11.4]{fri82} for the definition). Assume that the second étale homotopy group of $S$ is trivial e.g. if $S$ is the spectrum of a DVR or if it is an affine curve. It is possible to construct a pro-finite étale stack $\Pi_{X/S}\to \spec S$ with a morphism $X\to\Pi_{X/S}$ over $S$, called the \emph{relative étale fundamental gerbe} \cite[\S 8, \S 9]{bv15}, \cite[\S 1]{scfg} such that the isomorphism classes of points of the fibers are in natural bijection with the space of Galois sections of the étale fundamental group of the fiber. In short, $\Pi_{X/S}\to\spec S$ is a compact and coherent way of packing the spaces of Galois sections of the fibers.
